\section{Bốn bước và chỗ can thiệp}
\label{sec:ch10-bon-buoc}

\index{LIME!bốn bước của quy trình}%
Năm nhóm vấn đề cho biết một biến thể sinh ra để sửa cái gì, nhưng chưa cho biết
nó sửa ở đâu, và trục thứ hai của phân loại hỏi đúng câu ấy. Để hỏi được, khảo
sát chẻ LIME thành bốn bước và coi bốn bước đó là bốn cột của bảng phân loại;
một kỹ thuật được đánh dấu ở cột nào thì nó sửa bước
ấy~\autocite{p22whichlime}. Bốn bước như sau.

\index{sinh đặc trưng}%
\index{sinh mẫu}%
Bước một là \tn{sinh đặc trưng}{feature generation}: biến đầu vào thô thành tập
thành phần mà lời giải thích sẽ gọi tên, chẳng hạn bằng cách
\tn{phân đoạn}{segmentation} một tấm ảnh thành các super-pixel. Bước hai là
\tn{sinh mẫu}{sample generation}: dựng các mẫu đã nhiễu quanh điểm đang xét. Bước
ba là attribution đặc trưng: khớp một mô hình dễ đọc, thường là tuyến
tính, để xấp xỉ mô hình phức tạp trên vùng cục bộ ấy. Bước bốn là
\tn{biểu diễn lời giải thích}{explanation representation}: lấy trọng số của mô
hình thay thế mà trình ra \tn{mức quan trọng}{importance} của từng thành phần.

Bốn bước này chính là quy trình mà hình~\ref{fig:ch03-quy-trinh-lime} đã vẽ khi
chương~\ref{ch:lime} đọc bài báo gốc, chỉ gộp lại thô hơn một chút: bước dựng
lại đầu vào $z$ và bước gọi mô hình \tn{hộp đen}{black box} nằm gọn trong bước
hai của khảo sát,
còn phép khớp có trọng số nằm trong bước ba. Sự trùng khớp ấy không phải tình
cờ, vì cả hai đều đọc cùng một thủ tục từ cùng một bài báo. Điều đáng chú ý là
khảo sát dùng nó làm hệ tọa độ chứ không dùng làm hình minh họa: bốn bước trở
thành bốn chỗ mà một tác giả có thể đứng vào để sửa LIME.

Cách chẻ ấy nói được một điều mà danh sách tham số của
mục~\ref{sec:ch03-cac-tham-so-tu-do} chưa nói. Các tham số tự do không rải đều
trên quy trình; phép đối chiếu sau đây là của sách, vì khảo sát không nhắc tới
danh sách của chương~\ref{ch:lime}. Biểu diễn $x'$ thuộc bước một; phép dựng lại $z$, phân phối sinh
nhiễu và số mẫu $N$ thuộc bước hai; hàm khoảng cách $D$, độ rộng $\sigma$ và độ
dài $K$ thuộc bước ba; bước bốn không nhận lựa chọn nào trong bảy lựa chọn ấy. Vì vậy khi một
tác giả nói mình đã sửa được một nhóm vấn đề, chỗ đứng của họ trong bốn bước cho
biết họ đã đổi mệnh đề nào về $f$, và đó là thông tin mà tên của biến thể không
mang.

Khảo sát chia nhỏ hơn nữa bên trong từng bước, theo cách can thiệp. Ở bước một
có bốn nhóm: dựa trên phân đoạn, dựa trên phân cụm tập huấn luyện, dựa trên mức
quan trọng tính sẵn, và một nhóm mà khảo sát để ngỏ tên gọi, gồm những cách sinh
đặc trưng không theo ba lối trên, chẳng hạn sinh ra luật thay vì đặc trưng, hoặc
chạy song song hai kiểu dữ liệu. Ở bước hai có năm nhóm: chọn lọc mẫu, xấp xỉ,
dựa trên lân cận, dựa trên phân phối, và một nhóm để ngỏ nữa gồm những cách sinh
mẫu không theo bốn lối trên, chẳng hạn cách riêng cho một kiểu dữ liệu. Ở bước ba có bốn nhóm: sửa mô hình tuyến tính,
thay hẳn mô hình thay thế bằng cây, rừng hay đồ thị, đổi cách tính trọng số, và
đổi cách huấn luyện. Ở bước bốn có hai nhóm: mở rộng dạng lời giải thích, chẳng
hạn thành luật hay thành đồ thị, và làm lời giải thích tương tác
được~\autocite{p22whichlime}.

Danh sách con ấy đáng đọc vì nó cho thấy phạm vi của cái gọi là một biến thể của
LIME. Thay mô hình thay thế tuyến tính bằng một cây hồi quy đa đầu ra, hoặc bằng
một hệ lập trình logic quy nạp, tức một hệ học ra các luật logic từ ví dụ, thì
cái còn lại của LIME chỉ là ý tưởng nhiễu
quanh một điểm rồi khớp một mô hình dễ đọc. Bốn bước vì thế vừa là hệ tọa
độ vừa là lời cảnh báo: hai kỹ thuật cùng mang chữ LIME trong tên có thể không
còn chung một thủ tục nào ngoài cái khung bốn ô.
